%%%%%%%%%%%%%%%%%%%%%%%%%%%%%%%%%%%%%%%%%%%%%%%%%%%%%%%%%%%%%%%%%%%%%%%%%%%%%%%%
% Low-Budget Search-Guided Decision Making for Multi-Robot Pursuit--Evasion with Local Sensing
% Conference manuscript based on the IEEE conference template.
%
% Compile with pdfLaTeX / latexmk -pdf example.tex.
% The adjacent ieeeconf.cls is the unmodified PaperCept robotics template:
% https://ras.papercept.net/conferences/support/tex.php
% This is the IROS / RA-L initial-submission layout. RA-L accepted papers
% require a separate IEEEtran journal layout; do not mix the two styles.
% The current tables are retained as a record of the v5 benchmark while the
% corrected search implementation is being re-evaluated.
%%%%%%%%%%%%%%%%%%%%%%%%%%%%%%%%%%%%%%%%%%%%%%%%%%%%%%%%%%%%%%%%%%%%%%%%%%%%%%%%
%2345678901234567890123456789012345678901234567890123456789012345678901234567890
%        1         2         3         4         5         6         7         8

\documentclass[letterpaper,10pt,conference]{ieeeconf}

\IEEEoverridecommandlockouts
\overrideIEEEmargins

\usepackage[T1]{fontenc}
\usepackage{times}
\usepackage{graphicx}
\usepackage{amsmath}
\usepackage{amssymb}
\usepackage{bm}
\usepackage{microtype}
\usepackage{cite}
\usepackage{booktabs}
\usepackage{tabularx}
\usepackage{algorithm}
\usepackage{algpseudocode}

% Keep class-defined margins, headings, captions, and display spacing.
% Continuation lines in pseudocode align with the current block, without
% creating extra line numbers or losing the indentation of nested loops.
\algrenewcommand\algorithmicindent{1em}
\newcommand{\algcont}[1]{\Statex\hspace{\algorithmicindent}#1}
\newcolumntype{Y}{>{\raggedright\arraybackslash}X}

% Hardware sections are kept in the source until the physical evaluation is
% complete, but are hidden from the current manuscript.
\newif\ifshowhardware
\showhardwarefalse
\newcommand{\placeholderfig}[2]{%
  \fbox{\parbox[c][#1][c]{0.94\linewidth}{\centering\small\itshape #2}}}
% ------------------------------------------------------------------------------

\title{\LARGE\bfseries
Low-Budget Search-Guided Decision Making for Multi-Robot Pursuit--Evasion with Local Sensing
}

\author{Jincheng Shi$^{1}$ % <-this % stops a space
\thanks{$^{1}$Jincheng Shi is with the School of Automation,
        Beijing Institute of Technology, Beijing 100081, China.}%
}

\begin{document}

\maketitle
\thispagestyle{empty}
\pagestyle{empty}

%%%%%%%%%%%%%%%%%%%%%%%%%%%%%%%%%%%%%%%%%%%%%%%%%%%%%%%%%%%%%%%%%%%%%%%%%%%%%%%%
\begin{abstract}
Online search can improve adversarial decisions, but its cost grows with
the number of simulations. We study low-budget decision making for
continuous-space multi-robot pursuit--evasion by combining online search
with policies distilled from search experience. A team-alternating Monte
Carlo tree search labels states reached by the current policies; the
resulting labels train decentralized semantic-LiDAR policies and a value
model, which in turn guide candidate expansion and leaf evaluation during
deployment. The resulting controller can trade direct policy execution
against online lookahead as the available computation changes. Cross-play
experiments show that learned guidance and search budget affect pursuit
and evasion differently, motivating role-wise evaluation rather than a
single aggregate score. The formulation also defines the interface to a
kinodynamic execution layer for future robot experiments.

\end{abstract}

%%%%%%%%%%%%%%%%%%%%%%%%%%%%%%%%%%%%%%%%%%%%%%%%%%%%%%%%%%%%%%%%%%%%%%%%%%%%%%%%
\section{INTRODUCTION}

Multi-robot pursuit--evasion requires each team to coordinate while
responding to an opponent whose actions change the value of every joint
decision. Continuous actions, mutual elimination, and a finite decision
period make the problem difficult even in a bounded planar environment
\cite{isaacs1965differential,garcia2021multiple}.

Tree search evaluates candidate actions online and can adapt to a changing
opponent, but its cost grows with the simulation budget. A learned policy
amortizes computation offline and is fast at execution, but can lose the
lookahead needed for adversarial decisions. Search-guided learning methods
have combined these two forms of computation in multi-robot planning
\cite{riviere2021neural}. Here we ask how much of the benefit of lookahead
can be retained when learned guidance is used to allocate a smaller search
budget in a two-team game.

The two roles also make aggregate scores hard to interpret. A controller
that captures a weak evader need not resist a strong pursuer, and a long
episode does not by itself establish successful evasion. We therefore use
cross-play, in which every candidate controls both teams against every
other candidate, and report pursuit and evasion separately.

We formulate a bounded multi-pursuer multi-evader game with mutual
elimination at capture. Team-alternating MCTS supplies action labels on
states reached by the current policies. Per-agent policies are trained by
behavior cloning from these labels, while a value model predicts returns
for leaf evaluation. At deployment, policy proposals and value estimates
are mixed with ordinary progressive widening, so the same learned models
support direct execution and budgeted search. The policies use local
semantic-LiDAR observations; the current search and value modules use the
simulator state.

In a 2v2 simulation benchmark, the resulting controller provides an
intermediate operating point between direct policy execution and
unguided search. The cross-play results also expose a role-dependent
trade-off: increasing the search budget and adding learned guidance do
not affect pursuit and evasion in the same way.

The contributions of this work are:
\begin{enumerate}
  \item A team-alternating search-guided controller for continuous-action
        multi-robot pursuit--evasion, with vector backups and separate
        per-agent policy distillation.
  \item A deployment rule that combines local policy proposals and learned
        value estimates with progressive widening under a finite search
        budget.
  \item A two-sided cross-play evaluation that examines the budget--quality
        trade-off separately for pursuit and evasion.
\end{enumerate}

\section{RELATED WORK}

\subsection{Pursuit--Evasion Games}

PE games originate from Isaacs' differential games
\cite{isaacs1965differential}. Analytical and
semi-analytical solutions exist for structured cases, e.g. multi-player
reach-avoid games \cite{garcia2021multiple}, games with one superior
evader \cite{chen2016multiplayer}, and multi-pursuer interception
\cite{pierson2017intercepting}, but they rely on simplified dynamics and
full state information. Sensing limitations have been studied in
visibility-based formulations over polygonal environments
\cite{zhang2019tree}, which typically discretize space and
assume reactive rather than learned opponents. In contrast, we target
continuous space, continuous actions, and adversarial opponents that are
themselves produced by strong search.

\subsection{Tree Search for Adversarial Planning}

MCTS with UCT selection \cite{kocsis2006bandit,browne2012survey} is a
standard framework for online adversarial planning. Continuous actions are
handled with progressive widening \cite{couetoux2011continuous,
chaslot2008progressive}, while learned priors can bias expansion
\cite{rosin2011multi,auer2002finite}. Neural Tree Expansion combines
expert search with learned policy and value guidance for multi-robot
planning \cite{riviere2021neural}. We use the same broad teacher--student
idea in a mutual-elimination pursuit--evasion game, where the tree
alternates between teams and the policies are distilled independently for
each agent. The evaluation focuses on the effect of this design on the
budget--quality trade-off and on the difference between pursuit and
evasion.

\subsection{Learning-Based Multi-Agent Decision Making}

MARL methods such as MADDPG \cite{lowe2017multi} and MAPPO
\cite{yu2021surprising} learn policies through self-play
\cite{baker2020emergent} and can execute with little online computation.
Recent work has applied learned policies to pursuit--evasion under sensing
constraints \cite{gonultas2025pursuit,wang2024viper} and to multi-UAV
pursuit with online planning \cite{chen2025online}. These methods provide
useful fast baselines, but opponent-dependent evaluation remains essential
in adversarial tasks. Search-guided learning is also related to AlphaZero,
Expert Iteration, and behavior cloning from reachable states
\cite{silver2018general,anthony2017thinking,pomerleau1989alvinn,
ross2011reduction}. Our method combines these ingredients with a
team-alternating tree and decentralized action heads for the pursuit--
evasion setting.

%%%%%%%%%%%%%%%%%%%%%%%%%%%%%%%%%%%%%%%%%%%%%%%%%%%%%%%%%%%%%%%%%%%%%%%%%%%%%%%%
\section{PROBLEM FORMULATION}

\subsection{Two-Team Pursuit--Evasion Game}

We consider a multi-pursuer multi-evader game in a bounded planar region
$\Omega\subset\mathbb R^2$. The evader team
$\mathcal{E}=\{E_1,\ldots,E_{n_E}\}$ and the pursuer team
$\mathcal{P}=\{P_1,\ldots,P_{n_P}\}$ contain $n_E$ and $n_P$ robots,
respectively, with $N_r=n_E+n_P$ total agents. Equal-size $n$-versus-$n$
games are a special case. Index agents in evader-then-pursuer order.
The teams act at discrete decision times
$k$ with period $\Delta t$. Agent $i$ has position
$\mathbf{p}_{i,k}\in\mathbb{R}^2$ and alive flag
$z_{i,k}\in\{0,1\}$. The system state stacks positions, the elapsed time
$t_k$, and the alive flags:
\begin{equation}
\begin{aligned}
s_k={}&[\mathbf p_{1,k}^{T},\ldots,\mathbf p_{N_r,k}^{T},\\[-1mm]
      &\quad t_k,z_{1,k},\ldots,z_{N_r,k}]^{T}
      \in\mathbb R^{3N_r+1}.
\end{aligned}
\end{equation}
For the MINCO-enabled simulator we augment this decision state with the
two-dimensional velocity and acceleration of each robot,
\begin{equation}
\begin{aligned}
\bar s_k={}&[s_k^T,\mathbf v_{1,k}^T,\ldots,\mathbf v_{N_r,k}^T,\\[-1mm]
            &\quad\mathbf a_{1,k}^T,\ldots,\mathbf a_{N_r,k}^T]^T
            \in\mathbb R^{7N_r+1}.
\end{aligned}
\end{equation}
The decision state specifies the game variables. The simulator uses
the augmented representation in both execution modes; velocity and
acceleration components are zero in the direct-propagation benchmark.

Each agent commands a bounded 2-D displacement
\begin{equation}
a_{i,k}=\Delta\mathbf{p}_{i,k}\in\mathbb{R}^2,\qquad
\|a_{i,k}\|_\infty \le v_i^{\max}\Delta t ,
\end{equation}
and the joint action
\begin{equation}
\begin{aligned}
\mathbf a_k={}&[a_{1,k}^{T},\ldots,a_{N_r,k}^{T}]^{T}
                 \in\mathbb R^{2N_r},\\
s_{k+1}={}&F(s_k,\mathbf a_k).
\end{aligned}
\end{equation}
In direct propagation, alive agents advance by their bounded commands;
dead agents stay fixed. Commands within the action box may still produce
an invalid next state, which terminates the episode.

A capture occurs when a pursuer--evader pair comes within the capture
distance $d_c$:
\begin{equation}
\|\mathbf{p}_{P_i,k+1}-\mathbf{p}_{E_j,k+1}\|_2 < d_c .
\end{equation}
Multiple simultaneous captures are resolved by a nearest-first one-to-one
matching, and both members of a matched pair are deactivated
(\emph{mutual-elimination} rule). An episode terminates when either team
is fully inactive, the horizon $T$ is reached, or the state becomes
invalid:
\begin{equation}
\mathrm{Term}(s_{k+1}) =
\mathbb{I}\!\left[
 \begin{aligned}
 &N_E(s_{k+1})=0 \;\vee\; N_P(s_{k+1})=0\\[-1mm]
 &{}\vee\; t_{k+1}\ge T \;\vee\; \mathrm{invalid}(s_{k+1})
 \end{aligned}
\right],
\end{equation}
with $N_E(s_k)=\sum_{i\in\mathcal{E}}z_{i,k}$ and
$N_P(s_k)=\sum_{i\in\mathcal{P}}z_{i,k}$.

Let $c_k=N_E(s_k)-N_E(s_{k+1})$ count new captures. The shared
capture terms and horizon bonus are
\begin{equation}
\begin{aligned}
b_{P,k}&=w_c c_k(1-t_{k+1}/T),\\
b_{E,k}&=w_c c_k t_{k+1}/T
 +w_c N_E(s_{k+1})\mathbb I[t_k<T\le t_{k+1}],
\end{aligned}
\end{equation}
where $w_c$ is the capture-reward weight. Each active evader additionally
receives a distance-shaping term
\begin{equation}
g_{i,k}=w_d\,\operatorname{clip}_{[0,1]}
 \left(\frac{d_i(s_{k+1})-d_c}{d_{\max}-d_c}\right),
\end{equation}
where $d_i$ is the distance to the nearest active pursuer, $w_d=0.02$,
and $g_{i,k}=0$ when the evader is inactive or no active pursuer remains.
Thus $r_{i,k}=b_{E,k}+g_{i,k}$ for evaders and $r_{j,k}=b_{P,k}$ for
pursuers. If any member of a team violates the arena or obstacle
constraints, the reward of every member of that team is overwritten
with $-1$. Each team maximizes its own discounted return with
$\gamma=1$. Returns reported below sum rewards over the team members;
the distance term makes this game non-constant-sum.

\subsection{Observation Models}
\label{sec:observations}

The completed benchmark uses full simulator state for the search root
and value network. Per-agent distilled policies use the semantic LiDAR
encoding below; the MAPPO comparator uses normalized full state.
Consequently, full-state planning and local policy observations coexist: 
the fused solver is not an end-to-end partially observed planner.

In the \emph{local-perception} setting, agent $i$ receives
\begin{equation}
o_{i,k}=\phi_i(\bar s_k,\mathcal O)
=[\mathbf q_{i,k}^{T},\mathbf l_{i,k}^{T}]^{T},
\end{equation}
where $\mathcal{O}$ is a set of static rectangular obstacles,
$\mathbf{q}_{i,k}\in\mathbb{R}^8$ is a normalized ego encoding (alive
flag, position, velocity, acceleration, remaining time), and
$\mathbf{l}_{i,k}$ encodes a semantic LiDAR. The LiDAR casts $N_b$
uniformly spaced rays with fixed headings $\theta_m=2\pi m/N_b$ over
$N_c=4$ semantic channels (teammate, opponent, obstacle, boundary) with
range $d_{\max}$. Obstacles occlude dynamic targets: if a target and an
obstacle lie along ray $m$ at distances $d_{\mathrm{tar}}$ and
$d_{\mathrm{obs}}$, the effective reading is
\begin{equation}
d_{m,c} =
\begin{cases}
d_{\mathrm{tar}}, & d_{\mathrm{tar}} \le \min(d_{\mathrm{obs}}, d_{\max}),\\
d_{\max}, & \text{otherwise}.
\end{cases}
\end{equation}
The normalized matrix $L_{i,k}(m,c)=d_{m,c}/d_{\max}$ is flattened into
$\mathbf{l}_{i,k}$, giving a $264$-dimensional observation for
$N_b=64$; values near $0$ indicate nearby entities and $1$ no
detection.

\subsection{Objective and Deployment Constraints}

Each side seeks a per-agent policy
$\pi_i(o_{i,k})\mapsto a_{i,k}$ that maximizes its team return while
respecting a per-step computation budget $B$ (tree-search simulations)
and, on hardware, the latency budget of the decision period. We therefore
evaluate decision quality (capture rate, capture/survival time, return,
safety violations) jointly with online cost (per-step decision latency).

%%%%%%%%%%%%%%%%%%%%%%%%%%%%%%%%%%%%%%%%%%%%%%%%%%%%%%%%%%%%%%%%%%%%%%%%%%%%%%%%
\section{SEARCH-GUIDED POLICY DISTILLATION}

\subsection{Overview}

Fig.~\ref{fig:overview} summarizes the offline training and online
deployment stages.

\begin{figure*}[t]
\centering
\includegraphics[width=\textwidth]{figures/overview.pdf}
\caption{Overview of the search-guided controller. Offline, reachable
states are labeled by a team-alternating expert search and used to train
decentralized policy and value models. Online, local policy proposals and
learned leaf values guide a smaller search before the selected action is
executed by the dynamics layer.}
\label{fig:overview}
\end{figure*}

The framework closes a loop between three
modules---an expert search module, decentralized policies distilled
from it, and a value network estimating long-term team returns.
The standalone policies act from local encodings; the current fused
search operates on the simulator state at team level.
Offline, reachable states are annotated by a high-budget MCTS and
consumed by behavior cloning and value regression with time-weighted
replay; online, a small-budget search reuses the learned policy as a
candidate-action prior ($\beta_\pi$) and the value network as a leaf
evaluator ($\beta_v$). Algorithm~\ref{alg:sgpd} summarizes the loop.

\begin{algorithm}[t]
\caption{Search-Guided Policy Distillation (SGPD)}
\label{alg:sgpd}
\small
\begin{algorithmic}[1]
\Require Training rounds $L$; budgets $B_{\mathrm{train}}$,
\algcont{$B_{\mathrm{deploy}}$; guidance probabilities $\beta_\pi$, $\beta_v$}
\Ensure Per-agent policies $\hat\pi_{\theta_i}$; value model $V_\eta$
\For{$l=0$ \textbf{to} $L-1$}
  \State Roll out current policies to collect states
  \algcont{$\mathcal{S}_\pi=\{s^{(j)}\}$}
  \For{each $s^{(j)}\in\mathcal{S}_\pi$}
    \State Run team-alternating MCTS
    \algcont{with budget $B_{\mathrm{train}}$ (Sec.~\ref{sec:expert})}
    \State $c_j^*\gets\arg\max_c Q_c^{(\rho_0)}$
    \State $\mathbf{a}^{*(j)}\gets\mathbf{a}_{0\to c_j^*}$
    \State Roll out the current joint policy
    \algcont{from $s^{(j)}$ to termination}
    \State Compute return targets $\mathbf{v}^{*(j)}$ using
    \algcont{\eqref{eq:value_target}}
  \EndFor
  \State Update each $\theta_i$ using \eqref{eq:bc_loss}
  \algcont{with replay weights \eqref{eq:replay}}
  \State Update $\eta$ using \eqref{eq:value_loss}
\EndFor
\Statex \textbf{Online deployment (after training)}
\State At each step, run search with $B_{\mathrm{deploy}}$,
\algcont{policy proposals \eqref{eq:pw_guided} with probability $\beta_\pi$,}
\algcont{and learned leaf evaluation with probability $\beta_v$;}
\algcont{or execute $\hat\pi$ directly when $B_{\mathrm{deploy}}=0$}
\end{algorithmic}
\end{algorithm}

\subsection{Expert Action Generation by Team-Alternating MCTS}
\label{sec:expert}

A search node stores
\begin{equation}
\begin{aligned}
n={}&(s_n,\mathbf a_{p\to n},d_n,N_n,
       \mathbf W_n,\mathbf Q_n,\mathcal C_n),\\
\mathbf Q_n={}&\mathbf W_n/N_n .
\end{aligned}
\end{equation}
where $\mathbf W_n,\mathbf Q_n\in\mathbb R^{N_r}$ are cumulative and
mean value \emph{vectors} over the $N_r$ agents.
Edges carry full joint actions
\begin{equation}
\begin{aligned}
\mathbf a&=(\mathbf a^{\mathcal E},\mathbf a^{\mathcal P}),\\
\mathbf a^{\mathcal E}&=(a_{E_1},\ldots,a_{E_{n_E}}),\\
\mathbf a^{\mathcal P}&=(a_{P_1},\ldots,a_{P_{n_P}}).
\end{aligned}
\end{equation}
Every expansion advances the complete adversarial state
$s_c=F(s_n,\mathbf{a}_{n\to c})$. A root query is made for one team at
a time; the tree alternates the two team groups, and only the queried
team's slice of the selected joint action is executed.

Because the two teams have opposing interests, the evaluator alternates
by depth. With root evaluator $\rho_0$ and opponent $\bar{\rho}_0$,
\begin{equation}
\rho_n=
\begin{cases}
\rho_0, & d_n \bmod 2 = 0,\\
\bar{\rho}_0, & d_n \bmod 2 = 1,
\end{cases}
\label{eq:alternating}
\end{equation}
and branch selection compares only the scalar component
$Q_c^{(\rho_n)}$ of the value vector, so odd and even layers maximize the
opponent's and the root team's interests respectively.

Since the joint action space is continuous, full enumeration is
impossible. Each node maintains a finite candidate set
$\mathcal{A}(n)=\{\mathbf{a}_{n\to c}\mid c\in\mathcal{C}_n\}$ grown by
progressive widening,
\begin{equation}
M(n)=\left\lceil C_{\mathrm{pw}}\,N_n^{\alpha_{\mathrm{pw}}}\right\rceil ,
\label{eq:pw}
\end{equation}
with hyperparameters $(C_{\mathrm{pw}},\alpha_{\mathrm{pw}})$. New
candidates mix uniform sampling with the distilled policy prior. A shared
Bernoulli draw selects the proposal source for the joint candidate; in
the guided case, each active robot uses its own policy output:
\begin{equation}
\tilde a_i\sim
\begin{cases}
\hat\pi_{\theta_i}(\cdot\mid \phi_i(s_n)),&
\xi<\beta_\pi,\\
\mathcal U([a_i^{\min},a_i^{\max}]),&\text{otherwise}.
\end{cases}
\label{eq:pw_guided}
\end{equation}

Child selection follows a UCB rule from the current evaluator's
perspective,
\begin{equation}
\begin{aligned}
U(n,c;\rho_n)
={}&Q_c^{(\rho_n)}
 +C_{\mathrm{exp}}
\sqrt{\frac{\max(1,N_n)^{\alpha_{\mathrm{exp}}}}{N_c}} .
\end{aligned}
\label{eq:ucb}
\end{equation}
with exploration hyperparameters $(C_{\mathrm{exp}},\alpha_{\mathrm{exp}})$.

When a simulation reaches a leaf $s_L$ (depth limit or unexpanded node),
its value is estimated either by random rollout or by the value network:
with probability $\beta_v$ the leaf value is
$\hat{\mathbf{v}}_L=V_\eta(\psi(s_L))$, where $\psi$ denotes the
value-network input encoding (the full implementation state,
Sec.~\ref{sec:observations});
otherwise a rollout return is used. Estimates are discounted and backed up along the path, updating
$\mathbf{W}_n$, $\mathbf{Q}_n$ and $N_n$ component-wise. After the budget
is exhausted, the selected joint action is the root branch with highest
mean value for the queried evaluator,
\begin{equation}
c^*=\arg\max_{c\in\mathcal{C}_0} Q_c^{(\rho_0)},
\qquad
\mathbf{a}^*=\mathbf{a}_{0\to c^*} .
\label{eq:root_action}
\end{equation}

\subsection{Policy Distillation}

Directly fitting the joint action would couple $2N_r$ output dimensions
and scale poorly; we therefore maintain one policy network per agent.
Expert labels are collected on reachable states rather than uniform
samples: the current (or random) policy rolls out episodes, and the
visited states $\mathcal{S}_\pi=\{s^{(1)},\dots,s^{(N_\pi)}\}$ are
annotated by the expert. For state $s^{(j)}$ the expert joint action is
split per agent, $a_i^{*(j)}=\mathbf{a}^{*(j)}[\mathcal{I}_i]$, and
paired with the observation encoding
$x_i^{(j)}=\phi_i(s^{(j)})$, giving
\begin{equation}
\mathcal{D}_{\pi}^{(i)}
=
\left\{
\left(\phi_i(s^{(j)}),\mathbf a^{*(j)}[\mathcal I_i]\right)
\right\}_{j=1}^{N_\pi}.
\end{equation}

To mitigate the non-stationarity of self-generated data, samples enter a
replay buffer holding $T_{\mathrm r}$ rounds and are re-drawn with
linearly increasing time weights
\begin{equation}
p_t=\frac{t}{\sum_{\tau=1}^{T_{\mathrm r}}\tau},
\label{eq:replay}
\end{equation}
so recent expert knowledge is prioritized without discarding history.

Each policy is a compact three-layer 64-unit MLP per agent whose output
is squashed into the feasible action box by a scaled $\tanh$:
\begin{equation}
\hat{a}_i
=
\frac{a_i^{\max}+a_i^{\min}}{2}
+
\frac{a_i^{\max}-a_i^{\min}}{2}
\tanh\!\left(\mu_{\theta_i}(x_i)\right).
\end{equation}
Behavior cloning minimizes
\begin{equation}
\begin{aligned}
\mathcal{L}_{\pi}^{(i)}(\theta_i)
&=\frac{1}{N_\pi D_i}
\sum_{j=1}^{N_\pi}\sum_{d=1}^{D_i}
\left(\hat a_{i,d}^{(j)}-a_{i,d}^{*(j)}\right)^2 .
\end{aligned}
\label{eq:bc_loss}
\end{equation}
with AdamW \cite{loshchilov2019decoupled}, equivalently fitting the mean
under fixed variance. During standalone-policy evaluation, the raw mean
is perturbed by $0.1\epsilon$, $\epsilon\sim\mathcal N(0,I)$, before the
scaled $\tanh$. The policy serves both as a reactive baseline ($B=0$)
and as a candidate proposal in \eqref{eq:pw_guided}; MAPPO uses
deterministic inference.

\subsection{Value Network for Leaf Evaluation}

Deep random rollouts dominate the cost of online search. We replace them
with a value network $V_\eta$ trained on Monte Carlo return targets
collected from reachable states,
\begin{equation}
\mathbf{v}^{*(j)}
=\sum_{t=0}^{T_j-1}\gamma^t\,
\mathbf r_{t}^{\mathrm{norm},(j)} ,
\label{eq:value_target}
\end{equation}
by minimizing the element-wise error over all agents,
\begin{equation}
\begin{aligned}
\mathcal L_V(\eta)
&=\frac{1}{N_VN_r}
\sum_{j=1}^{N_V}\sum_{i=1}^{N_r}
\left(\hat v_i^{(j)}-v_i^{*(j)}\right)^2 .
\end{aligned}
\label{eq:value_loss}
\end{equation}
The return targets are Monte Carlo returns of the joint policy rolled out
from reachable states; the expert supervises actions only.
At deployment, leaves are evaluated by the network with probability
$\beta_v$ and by rollout otherwise (Sec.~\ref{sec:expert}), retaining rollout
evaluations alongside learned estimates.

\subsection{Kinodynamic Execution Layer}

Displacement commands cannot be executed instantaneously on hardware.
The simulator includes a
MINCO-based planner \cite{wang2022geometrically} that converts the high-level
command $\mathbf{u}_{i,k}$ into a smooth quintic trajectory
\begin{equation}
\mathbf p_i(t)=\mathbf C^T\bm\beta^{(0)}(t),\qquad
J=\int_0^{T_p}\|\mathbf p_i^{(3)}(t)\|_2^2\,dt .
\end{equation}
The coefficient state is propagated by
\begin{equation}
\mathbf C^+_{i,k}
=\mathbf F_{\mathrm{cl}}\mathbf C_{i,k}
+\mathbf G_{\mathrm{cl}}(\mathbf p^{\mathrm{cmd}}_{i,k})^T
\end{equation}
at the control period $\Delta t_c=\Delta t/N_s$. Raw targets are clipped
to the feasible region $[\underline{\mathbf{p}}_{i,k},
\overline{\mathbf{p}}_{i,k}]$ implied by velocity/acceleration limits and
the arena boundary before planning. Obstacle avoidance is left to the
high-level policy through its LiDAR observations, keeping the planner
closed-form; policy failures that cause collisions are penalized as
out-of-bounds terminations. The resulting simulator transition
    $s_{k+1}=F_{\mathrm{MINCO}}(s_k,\mathbf{u}_k)$ defines the execution
interface. The experiments below use direct propagation; the MINCO map is
included as the execution interface for later physical evaluation.

%%%%%%%%%%%%%%%%%%%%%%%%%%%%%%%%%%%%%%%%%%%%%%%%%%%%%%%%%%%%%%%%%%%%%%%%%%%%%%%%
\section{EXPERIMENTS}

\subsection{Protocol and Evaluation Questions}

The evaluation asks two questions. First, how does search-guided policy
execution compare with direct policy execution and unguided search at the
tested operating points? Second, do the same comparisons hold for pursuit
and evasion? We answer these questions with a complete cross-play matrix.
Matched-budget and hardware measurements are treated separately from this
comparison.

The completed v5 benchmark instantiates the formulation with
$n_E=n_P=2$ (2v2), giving a 13-D decision state, a 29-D augmented state,
and an eight-dimensional joint action. Team size is an evaluation
configuration rather than a defining restriction of the framework.
The v5 protocol uses a $[-10,10]^2$ arena, an 80-step horizon,
$\Delta t=1$, capture distance $d_c=1$, capture-reward weight $w_c=0.5$,
and mutual elimination.
Actions satisfy box bounds with evader speed parameter 1.5 and pursuer
parameter 1.0. Initial positions are sampled in $[-9,9]$ in $x$, with
evaders in $y\in[0,9]$ and pursuers in $y\in[-9,0]$, rejecting invalid
or insufficiently separated configurations. The benchmark contains no
obstacles and uses direct propagation, retaining the 29-D state with
zero velocity and acceleration components.

All methods face the same initial-state distribution. Ordinary matchups
use seeds 1000--1004 with 200 episodes each; matchups involving the
50,000-simulation reference use seed 1000 and 200 episodes. These are
evaluation seeds, not independent training seeds. The five-by-five
cross-play contains 17,800 episodes. We report double capture, mean
number captured where relevant, unconditional episode duration,
team-specific return, and invalid termination. An episode that avoids
double capture may still lose one evader; duration is not individual
robot survival. The saved summaries use $1.96s/\sqrt n$ intervals. An
observed 100\% rate has zero sample variance but does not imply certainty.
The evaluated distilled model is checkpoint iteration 79 from the
\texttt{18\_lidar\_2obstacle} family; the baseline uses
\texttt{mappo\_50k.pt}. Policy inputs are 264-D semantic LiDAR encodings,
while search and value evaluation receive full simulator state. Thus the
current experiment evaluates local policy inputs together with a
full-state search root; it does not evaluate a fully partially observed
search planner.

\begin{table}[t]
\centering\footnotesize
\caption{Completed cross-play configuration. Execution-layer settings
are separate from this direct-propagation benchmark.}
\label{tab:params}
\begin{tabularx}{\columnwidth}{@{}Yl@{}}\toprule
Parameter & Value \\ \midrule
Evader / pursuer team size & $n_E=n_P=2$ \\
Arena / horizon & $[-10,10]^2$ / 80 steps \\
Evader / pursuer speed parameter & 1.5 / 1.0 \\
Capture radius / discount & 1.0 / 1.0 \\
Search depth & 100 \\
SO / S+L / LM simulations & 3000 / 1000 / 50000 \\
$(C_{\mathrm{pw}},\alpha_{\mathrm{pw}})$ & $(2.0,0.5)$ \\
$(C_{\mathrm{exp}},\alpha_{\mathrm{exp}})$ & $(1.0,0.25)$ \\
$(\beta_\pi,\beta_v)$ for S+L & $(0.5,0.5)$ \\
Policy / value input dimension & 264 / 29 \\
Policy LiDAR / obstacles & 64 rays / 0 \\
\bottomrule\end{tabularx}
\end{table}

\subsection{Compared Decision Makers}

The completed benchmark compares five decision makers. Policy Only and
the policy-guided expansion use local semantic-LiDAR policy encodings;
the current C++ search root and value evaluator use simulator state.
\begin{itemize}
\item \textbf{Policy Only}: distilled policies without online search.
\item \textbf{Search Only}: unguided team-alternating MCTS with uniform
candidate sampling, random-rollout leaf evaluation, and $B=3000$.
\item \textbf{Search + Learning}: the same search family with learned
candidate proposals and leaf evaluation, $(\beta_\pi,\beta_v)=(0.5,0.5)$,
and $B=1000$.
\item \textbf{MAPPO}: a self-play-trained shared-actor comparator
\cite{yu2021surprising}, using normalized full-state inputs. The current
implementation lacks an agent identifier or ego-conditioned input:
robots on the same team receive identical inputs to the same actor and
therefore execute identical deterministic normalized actions. Its
results describe this implementation, not a definitive MAPPO comparison.
\item \textbf{Large-budget MCTS}: unguided search at $B=50000$,
used as an empirical high-computation reference rather than an optimal
solution or a theoretical performance bound.
\end{itemize}

\subsection{Cross-Play Reveals Role-Dependent Benefits}

Table~\ref{tab:matrix} reports all pairings, rather than selecting a
single favorable opponent. For more detailed analysis, we use the S+L
evader, whose column distinguishes pursuit performance, and the
50,000-simulation pursuer, which achieves the highest double-capture
rate against that evader among the evaluated methods
(Table~\ref{tab:twosided}). Reference selection is based on this candidate
set; it is not a claim of universal opponent strength.

\begin{table*}[t]
\centering\footnotesize
\caption{Complete v5 cross-play: double-capture rate (\%). Rows are
pursuers and columns are evaders. PO: Policy Only; SO: Search Only;
S+L: Search + Learning; LM: Large-budget MCTS. Ordinary matchups have
1,000 episodes; any matchup involving LM has 200. $^*$Current MAPPO
implementation; see limitations in the text.}
\label{tab:matrix}
\begin{tabular*}{\textwidth}{@{\extracolsep{\fill}}lrrrrr@{}}\toprule
Pursuer $\backslash$ Evader & PO & SO & S+L & MAPPO$^*$ & LM \\ \midrule
Policy Only ($B=0$) & 100.0 & 100.0 & 31.0 & 98.3 & 100.0 \\
Search Only ($B=3000$) & 100.0 & 99.9 & 60.9 & 98.6 & 100.0 \\
Search + Learning ($B=1000$) & 100.0 & 100.0 & 49.9 & 98.5 & 100.0 \\
MAPPO$^*$ & 4.6 & 46.2 & 0.1 & 23.0 & 50.0 \\
Large-budget MCTS ($B=50000$) & 100.0 & 100.0 & 76.0 & 98.0 & 100.0 \\
\bottomrule\end{tabular*}
\end{table*}

\begin{table*}[t]
\centering\footnotesize
\caption{Two-sided reference evaluation, mean $\pm$ episode-level
95\% normal-approximation CI from the saved v5 summaries. Duration is
unconditional episode length. Invalid denotes termination by a
collision or boundary violation on either side.}
\label{tab:twosided}
\begin{tabular*}{\textwidth}{@{\extracolsep{\fill}}lrrrr@{}}\toprule
\multicolumn{5}{l}{\textbf{(a) Pursuit vs. S+L evader}}\\
Method & Double cap. (\%) & Duration (steps) & Pursuer return & Invalid (\%) \\ \midrule
Policy Only ($B=0$) & $31.0\pm2.9$ & $74.5\pm0.7$ & $0.428\pm0.027$ & 0.0 \\
Search Only ($B=3000$) & $60.9\pm3.0$ & $65.0\pm1.0$ & $0.791\pm0.022$ & 0.0 \\
Search + Learning ($B=1000$) & $49.9\pm3.1$ & $64.3\pm1.2$ & $0.751\pm0.026$ & 0.0 \\
MAPPO$^*$ & $0.1\pm0.2$ & $79.9\pm0.1$ & $0.139\pm0.019$ & 0.1 \\
Large-budget MCTS ($B=50000$) & $76.0\pm5.9$ & $56.6\pm2.8$ & $1.001\pm0.048$ & 0.0 \\
\midrule
\multicolumn{5}{l}{\textbf{(b) Evasion vs. 50,000-simulation pursuer}}\\
Method & Double caught (\%) & Duration (steps) & Evader return & Invalid (\%) \\ \midrule
Policy Only ($B=0$) & $100.0\pm0.0$ & $17.2\pm0.9$ & $0.421\pm0.018$ & 0.0 \\
Search Only ($B=3000$) & $100.0\pm0.0$ & $13.6\pm0.7$ & $0.324\pm0.013$ & 0.0 \\
Search + Learning ($B=1000$) & $76.0\pm5.9$ & $56.6\pm2.8$ & $1.197\pm0.055$ & 0.0 \\
MAPPO$^*$ & $98.0\pm1.9$ & $16.1\pm0.9$ & $0.310\pm0.046$ & 2.0 \\
Large-budget MCTS ($B=50000$) & $100.0\pm0.0$ & $11.4\pm0.6$ & $0.267\pm0.012$ & 0.0 \\
\bottomrule\end{tabular*}
\end{table*}

\paragraph*{Cross-play results}
The matrix shows a clear role dependence in the tested operating points.
Against the S+L evader, pursuit improves as the search budget increases:
the policy-only controller is weakest, the 1,000-simulation fused search
is intermediate, and the larger unguided searches perform better. Against
the high-budget pursuer, the fused evader lasts substantially longer than
the policy-only and unguided-search evaders. These observations motivate
separate pursuit and evasion analyses; they do not establish a universal
ordering of the methods.

The MAPPO row is included as an implementation baseline. Its shared actor
does not receive an agent identifier, so same-team agents receive the same
deterministic action for identical inputs. The corresponding results are
useful for describing this implementation, but are not evidence about
MAPPO in general. Invalid terminations are retained in the reported
statistics.

The tables report the v5 benchmark. Because the search implementation used
for this run was subsequently corrected, the quantitative ordering should
be regarded as provisional pending a rerun with the corrected implementation.

% The following sections are retained as an experiment worklist.  They are
% enabled only after the corrected benchmark and hardware measurements are
% available.
\iffalse
\subsection{Matched-Budget and Component Studies}

The completed comparison changes both search budget and guidance.
The v5 result establishes a low-simulation operating point, but not a
matched-time advantage. To isolate budget from guidance, the next evaluation fixes the checkpoint,
opponent, initial conditions, search depth, and reward. SO and S+L will
be evaluated at $B\in\{200,500,1000,3000,5000\}$ on both roles.
The resulting curves will report capture performance and episode
duration against simulation count and independently measured decision
time. Earlier thesis-scale sweeps used different opponents and execution
protocols; we exclude their numbers from the main comparison.
Run the matched-budget sweep under the v5 protocol. Do not
claim matched-time superiority or a speedup from simulation ratios alone.

At $B=1000$, a factorial component study separates policy-based proposals
from value-based evaluation (Table~\ref{tab:beta}). The opponent remains
fixed at $(0.5,0.5)$, while the tested side uses $(0,0)$, $(0.5,0)$,
$(0,0.5)$, or $(0.5,0.5)$. Pursuit and evasion are evaluated separately.
The same study is repeated with two occluding obstacles and the execution
layer enabled, with dynamics and model-training provenance specified
for each scenario. A $(1,1)$ setting can additionally test dependence
on rollout fallback.

\begin{table}[t]\centering\footnotesize
\caption{Planned component ablation at $B=1000$. Pending cells are not
experimental results. Repeat on both the direct-propagation and
obstacle/execution scenarios.}
\label{tab:beta}
\begin{tabular}{@{}ccccc@{}}\toprule
$\beta_\pi$ & $\beta_v$ & Pursuit dbl. (\%) & Duration & Invalid (\%)\\
\midrule
0 & 0 & -- & -- & --\\
0.5 & 0 & -- & -- & --\\
0 & 0.5 & -- & -- & --\\
0.5 & 0.5 & -- & -- & --\\
\bottomrule\end{tabular}
\end{table}
Add the corresponding evasion table: double-caught rate,
episode duration, evader return, and invalid rate.

To examine the rollout-mismatch hypothesis, fixed candidate actions at
representative states will be ranked by random-rollout estimates,
learned values, and returns against the reference pursuer. Comparing
ranking consistency and variance, together with random versus
search-based rollout opponents, will test whether evaluation quality
explains the role difference. The mechanism claim will be retained only
if this controlled evidence supports it.

\subsection{Observation and Execution Validation}

The policy's semantic rays are subject to obstacle occlusion. The
current search and value model still receive full simulator state, so
local sensing is presently a policy interface rather than a complete
partial-observation search solution. The correct next step is an
estimated-state or belief-state root for both search and value evaluation.
Accordingly, an obstacle test can study robustness of local action
proposals without establishing a belief-space planner. Before claiming
local-only fused deployment, the root-state estimator and the treatment
of unobserved opponents must be specified and evaluated. A
full-state-versus-estimated-state comparison should distinguish planning
information from the policy observation interface.
Report obstacle/execution simulation on the same initial-state
sets for Policy Only, SO, and S+L. Include collisions, boundary exits,
trajectory error, and capture outcomes; verify Python/C++ transition
consistency. Earlier isolated obstacle numbers are not pooled with v5.

The v5 run records each team's solver-call time separately, correcting
joint-action timing. These measurements were collected concurrently
across workers, however, and are not controlled onboard latency.
A dedicated measurement will fix hardware, device, threads, warm-up,
and a common state set. The benchmark records solver-call samples after
a configurable warm-up and reports the exact pooled quantities
$\bar\tau_r=\sum_{q\in\mathcal Q_r}\tau_q/|\mathcal Q_r|$ for each
role $r\in\{\mathrm{evader},\mathrm{pursuer},\mathrm{joint}\}$,
together with per-episode p50/p95 diagnostics. The pooled mean is the
reported average single-step decision time; it excludes environment
transition and reward computation, and the joint value includes both
team solver calls executed sequentially. The mean of episode-level 95th
percentiles is not the pooled per-call 95th percentile.
Fill the dedicated latency table and include perception,
communication, search, and trajectory generation in end-to-end timing.

\ifshowhardware
\subsection{Physical-Robot Interface (Results Reserved)}

This subsection is an interface for the future physical experiment; it
contains no hardware result. The experiment will test whether the
simulation behavior survives sensing, communication, and execution
errors. The intended comparison is Policy Only versus S+L on the same
platforms, opponents, layouts, and initial configurations; a feasible
small-budget SO setting may be added as a search-only reference. The
high-level decision period must be distinguished from the low-level
trajectory-control period. For any fused deployment, state acquisition
and unavailable opponent states will be documented explicitly.

Table~\ref{tab:hardware_interface} defines the contract between the
decision module and the robot stack. Each field is deliberately left
parameterized so that the eventual hardware measurements can be inserted
without changing the algorithm description.

\begin{table}[t]
\centering\footnotesize
\caption{Physical-experiment interface. Entries to be completed from the
filled from the robot platform and trial log; no value is assumed here.}
\label{tab:hardware_interface}
\begin{tabular}{@{}p{0.25\linewidth}p{0.68\linewidth}@{}}
\toprule
Interface item & Definition / field to be specified \\
\midrule
Input to decision module & $o_{i,k}$: semantic LiDAR, robot state,
timestamp; sensor and localization fields \\
Opponent information & measured, communicated, or estimated
state and delay \\
Output & $\mathbf u_{i,k}$: bounded displacement target and validity flag \\
Execution layer & MINCO target-to-trajectory map
 $\mathbf u_{i,k}\mapsto\mathbf p_i(t)$ \\
Timing & decision period, control period, compute device,
mean/p50/p95, deadline misses \\
Safety log & collision, boundary, tracking error, emergency stop \\
Trial protocol & robot count, layouts, initial states, trials,
success criteria \\
\bottomrule
\end{tabular}
\end{table}

\begin{figure}[t]\centering
\placeholderfig{2.4cm}{Pending hardware figure: arena and robot teams;
perception/state-estimation interface; high-level decision and MINCO
control timing; representative capture and escape trajectories.}
\caption{Physical validation layout and execution interface to be
completed after measurements.}
\label{fig:hardware}
\end{figure}

The completed subsection will report team sizes, trial counts, captured
fraction, full-team capture or avoidance of full-team capture,
collision and boundary violations,
trajectory-tracking error, and end-to-end latency with deadline misses.
Open and obstacle layouts will be evaluated separately. Failure cases
will identify whether the cause is state estimation, decision selection,
or execution. If only Policy Only is deployed, the hardware claim will
be limited to distilled-policy execution rather than validation of
online fusion.
Add robot platform, sensing/localization, onboard computer,
control and decision periods, trial protocol, outcomes, and video. No
physical result is claimed in this manuscript version.
\fi

\fi

\subsection{Scope of the Current Evidence}

The benchmark evaluates one 2v2 configuration and one distilled
checkpoint family. The search root and value model use simulator state,
whereas the policies use semantic LiDAR, so the experiment does not yet
represent a fully partially observed fused planner. The MAPPO comparator
also has a restricted actor input. These conditions delimit the scope of
the cross-play evidence and motivate the matched-budget, component,
state-estimation, and hardware evaluations retained in the source below.


\section{CONCLUSION}

We presented a search-guided controller for continuous-action
multi-robot pursuit--evasion. Team-alternating MCTS provides supervision
for decentralized policies, and the learned policies and value model guide
subsequent low-budget search. In the cross-play benchmark, the
controller occupies an intermediate operating point between direct policy
execution and larger unguided searches, while pursuit and evasion exhibit
different responses to search and learned guidance.

The evidence is limited to the stated 2v2 simulation protocol and a
full-state search root. Re-evaluation with the corrected implementation,
matched budgets, component ablations, and physical execution will determine
how well the observed trade-off transfers beyond this setting.

%%%%%%%%%%%%%%%%%%%%%%%%%%%%%%%%%%%%%%%%%%%%%%%%%%%%%%%%%%%%%%%%%%%%%%%%%%%%%%%%

\begin{thebibliography}{99}

\bibitem{isaacs1965differential}
R.~Isaacs, \emph{Differential Games: A Mathematical Theory with
Applications to Warfare and Pursuit, Control and Optimization}. New York,
NY, USA: Wiley, 1965.

\bibitem{garcia2021multiple}
E.~Garcia, D.~W. Casbeer, A.~von Moll, and M.~Pachter, ``Multiple
pursuer multiple evader differential games,'' \emph{IEEE Trans. Autom.
Control}, vol.~66, no.~5, pp.~2345--2350, 2021.

\bibitem{chen2016multiplayer}
J.~Chen, W.~Zha, Z.~H. Peng, et al., ``Multi-player pursuit--evasion
games with one superior evader,'' \emph{Automatica}, vol.~71, pp.~24--32,
2016.

\bibitem{pierson2017intercepting}
A.~Pierson, Z.~J. Wang, and M.~Schwager, ``Intercepting rogue robots: An
algorithm for capturing multiple evaders with multiple pursuers,''
\emph{IEEE Robot. Autom. Lett.}, vol.~2, no.~2, pp.~530--537, 2017.

\bibitem{zhang2019tree}
Z.~Zhang, Y.~Sung, L.~Zhou, et al., ``Tree search techniques for
minimizing detectability and maximizing visibility,'' in \emph{Proc.
IEEE Int. Conf. Robot. Autom. (ICRA)}, Montreal, QC, Canada, 2019.

\bibitem{kocsis2006bandit}
L.~Kocsis and C.~Szepesv\'ari, ``Bandit based Monte-Carlo planning,'' in
\emph{Proc. Eur. Conf. Mach. Learn. (ECML)}, 2006, pp.~282--293.

\bibitem{browne2012survey}
C.~B. Browne et al., ``A survey of Monte Carlo tree search methods,''
\emph{IEEE Trans. Comput. Intell. AI Games}, vol.~4, no.~1, pp.~1--43,
2012.

\bibitem{couetoux2011continuous}
A.~Cou\"etoux, J.-B. Hoock, N.~Sokolovska, O.~Teytaud, and N.~Bonnard,
``Continuous upper confidence trees,'' in \emph{Proc. LION}, 2011,
pp.~433--445.

\bibitem{chaslot2008progressive}
G.~M. J.-B. Chaslot, M.~H. M. Winands, J.~W. H. M. Uiterwijk, H.~J. van
den Herik, and B.~Bouzy, ``Progressive strategies for Monte-Carlo tree
search,'' \emph{New Math. Nat. Comput.}, vol.~4, no.~3, pp.~343--357,
2008.

\bibitem{rosin2011multi}
C.~D. Rosin, ``Multi-armed bandits with episode context,'' \emph{Ann.
Math. Artif. Intell.}, vol.~61, no.~3, pp.~203--230, 2011.

\bibitem{riviere2021neural}
B.~Riviere, W.~H\"onig, M.~Anderson, and S.-J. Chung, ``Neural tree
expansion for multi-robot planning in non-cooperative environments,''
\emph{IEEE Robot. Autom. Lett.}, vol.~6, no.~4, pp.~6868--6875, 2021.

\bibitem{huang2022more}
B.~Huang, T.~Guo, A.~Boularias, and J.~Yu, ``Interleaving Monte Carlo
tree search and self-supervised learning for object retrieval in clutter,''
arXiv:2202.01426, 2022.

\bibitem{silver2018general}
D.~Silver et al., ``A general reinforcement learning algorithm that
masters chess, shogi, and Go through self-play,'' \emph{Science},
vol.~362, no.~6419, pp.~1140--1144, 2018.

\bibitem{anthony2017thinking}
T.~Anthony, Z.~Tian, and D.~Barber, ``Thinking fast and slow with deep
learning and tree search,'' in \emph{Proc. Adv. Neural Inf. Process.
Syst. (NeurIPS)}, 2017, pp.~5360--5370.

\bibitem{yu2021surprising}
C.~Yu et al., ``The surprising effectiveness of PPO in cooperative
multi-agent games,'' in \emph{Proc. NeurIPS Datasets and Benchmarks},
2021.

\bibitem{wang2022geometrically}
Z.~Wang, X.~Zhou, C.~Xu, and F.~Gao, ``Geometrically constrained
trajectory optimization for multicopters,'' \emph{IEEE Trans. Robot.},
vol.~38, no.~5, pp.~3259--3278, 2022.

\bibitem{loshchilov2019decoupled}
I.~Loshchilov and F.~Hutter, ``Decoupled weight decay regularization,''
in \emph{Proc. Int. Conf. Learn. Represent. (ICLR)}, 2019.

\bibitem{auer2002finite}
P.~Auer, N.~Cesa-Bianchi, and P.~Fischer, ``Finite-time analysis of the
multiarmed bandit problem,'' \emph{Mach. Learn.}, vol.~47, no.~2--3,
pp.~235--256, 2002.

\bibitem{ross2011reduction}
S.~Ross, G.~J. Gordon, and J.~A. Bagnell, ``A reduction of imitation
learning and structured prediction to no-regret online learning,'' in
\emph{Proc. Int. Conf. Artif. Intell. Stat. (AISTATS)}, 2011,
pp.~627--635.

\bibitem{pomerleau1989alvinn}
D.~A. Pomerleau, ``ALVINN: An autonomous land vehicle in a neural
network,'' in \emph{Proc. Adv. Neural Inf. Process. Syst. (NeurIPS)},
1989, pp.~305--313.

\bibitem{lowe2017multi}
R.~Lowe, Y.~Wu, A.~Tamar, J.~Harb, P.~Abbeel, and I.~Mordatch,
``Multi-agent actor-critic for mixed cooperative-competitive
environments,'' in \emph{Proc. Adv. Neural Inf. Process. Syst.
(NeurIPS)}, 2017, pp.~6379--6390.

\bibitem{schrittwieser2020mastering}
J.~Schrittwieser, I.~Antonoglou, T.~Hubert, et al., ``Mastering Atari,
Go, chess and shogi by planning with a learned model,'' \emph{Nature},
vol.~588, no.~7839, pp.~604--609, 2020.

\bibitem{baker2020emergent}
B.~Baker, I.~Kanitscheider, T.~Markov, et al., ``Emergent tool use from
multi-agent autocurricula,'' in \emph{Proc. Int. Conf. Learn. Represent.
(ICLR)}, 2020.

\bibitem{gonultas2025pursuit}
B.~M. Gonultas and V.~Isler, ``Pursuit-evasion for car-like robots with
sensor constraints,'' in \emph{Proc. IEEE/RSJ Int. Conf. Intell. Robots
Syst. (IROS)}, 2025.

\bibitem{wang2024viper}
Y.~Wang, Y.~Cao, J.~Chiun, S.~Koley, M.~Pham, and G.~Sartoretti,
``ViPER: Visibility-based pursuit-evasion via reinforcement learning,''
in \emph{Proc. Conf. Robot Learn. (CoRL)}, ser. PMLR, vol.~270, 2024,
pp.~4188--4200.

\bibitem{chen2025online}
J.~Chen, C.~Yu, G.~Li, et al., ``Online planning for multi-UAV
pursuit-evasion in unknown environments using deep reinforcement
learning,'' \emph{IEEE Robot. Autom. Lett.}, 2025.

\end{thebibliography}

\end{document}
